\section{Oficiálne zadanie}
\label{sec:zadanie}

\paragraph{Názov:} Inteligentný systém pre automatizované hodnotenie
bezpečnosti a~governance AI systémov

\paragraph{Cieľ práce:} Tímový projekt v~spolupráci so spoločnosťou Alison
Slovakia s.r.o. Cieľom je navrhnúť a~implementovať systém, ktorý dokáže
automatizovane zhromažďovať a~vyhodnocovať informácie o~AI modeli alebo AI
systéme, vytvárať jeho bezpečnostný a~governance profil a~identifikovať
riziká, nedostatky a~chýbajúce informácie na základe dostupných dôkazov.

\paragraph{Anotácia:} Informácie potrebné na posúdenie AI systému sú často
rozdelené medzi model cards, technickú dokumentáciu, informácie o~použitých
dátach, evaluačné výsledky, bezpečnostné deklarácie a~ďalšie zdroje. Ich
manuálne spracovanie je časovo náročné a~výsledné hodnotenie môže výrazne
závisieť od skúseností hodnotiteľa. Projekt sa zameria na vytvorenie
inteligentného systému, ktorý bude tieto informácie automatizovane získavať,
prepájať a~využívať pri hodnotení bezpečnosti a~governance AI riešení. Systém
môže vytvárať štruktúrovaný profil AI modelu alebo systému, automaticky
dopĺňať model card, identifikovať relevantné riziká, vytvárať threat model
alebo vykonávať ďalšie formy bezpečnostného a~governance hodnotenia. Ak pre
rozhodnutie nebudú dostupné dostatočné informácie, systém bude schopný
identifikovať chýbajúce informácie, generovať doplňujúce otázky alebo
požiadavky na ďalšie zdroje. Zistenia budú mapované na vybrané bezpečnostné
a~governance požiadavky a~spätne dohľadateľné k~dôkazom, z~ktorých vznikli.
Výsledkom bude prototyp adaptívneho hodnotiaceho systému, ktorý dokáže
postupne vytvárať a~spresňovať bezpečnostný a~governance profil AI riešenia.

\paragraph{Úlohy:}
\begin{enumerate}
    \item Navrhnúť jednotný model informácií o~AI riešení zahŕňajúci vybrané
    oblasti: model card, účel a~obmedzenia použitia, informácie o~dátach,
    evaluačné výsledky, bezpečnostné vlastnosti, ďalšie relevantné
    charakteristiky hodnoteného riešenia.
    \item Implementovať automatizovaný zber, extrakciu a~prepájanie informácií
    z~model cards, technickej dokumentácie a~ďalších dostupných zdrojov.
    \item Implementovať vybrané automatizované hodnotiace mechanizmy: tvorbu
    alebo dopĺňanie model card, threat modeling, identifikáciu bezpečnostných
    rizík, kontrolu splnenia definovaných požiadaviek.
    \item Implementovať adaptívne hodnotenie schopné identifikovať chýbajúce
    informácie, generovať doplňujúce otázky a~pracovať s~mierou istoty
    jednotlivých zistení.
    \item Vyhodnotiť systém na vybraných AI modeloch alebo systémoch z~pohľadu
    správnosti a~úplnosti získaných informácií, kvality vytvorených hodnotení,
    identifikovaných rizík a~dohľadateľnosti zistení k~zdrojovým dôkazom.
\end{enumerate}
